\section{Experimental Setup}
\label{sec:setup}

\subsection{Benchmark instances}
\label{ssec:instances}

The experiments are conducted on instances from the DIMACS graph
coloring benchmark suite. These instances are widely used for
evaluating graph coloring algorithms and cover several graph families.
The DSJC instances are random graphs, where the suffix indicates the
edge density. The \texttt{le450} instances are Leighton graphs with
structured coloring constraints. The flat instances are quasi-random
graphs generated with a prescribed target number of colors and
relatively balanced vertex degrees. Table~\ref{tab:instances} lists
the benchmark instances used in the experiments. The second column
reports the pair $\chi/k^\star$, where $\chi$ denotes the chromatic
number when it is known (``?'' indicates that it is unknown) and
$k^\star$ is the number of colors in the best known solution.

\begin{table}[t]
\centering
\caption{DIMACS benchmark instances used in the experiments. Notice the size group is based on the number of vertices: Small for
$n\leq 300$, Medium for $300<n\leq 500$, and Large for $n>500$.}
\label{tab:instances}
\small
\setlength{\tabcolsep}{5pt}
\begin{tabular}{@{}lccrl@{}}
\toprule
Instance & $\chi/k^\star$ & Size & $n$ & Family \\
\midrule
\texttt{DSJC250.5}        & ?/28   & Small  & 250  & Random, density 0.5 \\
\texttt{DSJC500.1}        & ?/12   & Medium & 500  & Random, density 0.1 \\
\texttt{DSJC500.5}        & ?/47   & Medium & 500  & Random, density 0.5 \\
\texttt{DSJC500.9}        & ?/126  & Medium & 500  & Random, density 0.9 \\
\texttt{flat300\_26\_0}   & 26/26  & Small  & 300  & Flat graph \\
\texttt{flat300\_28\_0}   & 28/28  & Small  & 300  & Flat graph \\
\texttt{flat1000\_76\_0}  & 76/81  & Large  & 1000 & Flat graph \\
\texttt{le450\_15c}       & 15/15  & Medium & 450  & Leighton graph \\
\texttt{le450\_15d}       & 15/15  & Medium & 450  & Leighton graph \\
\texttt{le450\_25c}       & 25/25  & Medium & 450  & Leighton graph \\
\texttt{le450\_25d}       & 25/25  & Medium & 450  & Leighton graph \\
\bottomrule
\end{tabular}
\end{table}

\subsection{Experimental settings and evaluation metrics}
\label{ssec:exp-settings}

Detailed parameter settings for RLTabucol, RLACO, and PPO are reported
in Appendix~A. Training is performed only on Erd\H{o}s--R\'enyi random
graphs using a T4 GPU; no DIMACS instance is used during training.
Testing is conducted on CPU, and all reported times are wall-clock times.
Unless stated otherwise, each experimental configuration is evaluated
with ten independent runs.

For the RLTabucol-related framework experiment, RLTabucol, RLTabucol$_0$, and
Pure Tabucol are evaluated on the same tested color budgets with a
five-hour ($18{,}000$\,s) time limit per run. The five-hour time limit follows the source RLTabucol evaluation protocol
and is consistent with long run-time budgets used in previous DIMACS
graph-coloring studies~\cite{frijters2023rltabucol,zhou2018plscoll}. The primary metric
is $k_{\mathrm{best}}$, the smallest tested color budget for which a
legal coloring ($f(c)=0$) is found. At the same $k$, success count and
successful-run time are used as secondary measures. Success denotes
the number of zero-conflict runs out of ten; time is reported as the
best and mean wall-clock time over successful runs only. The actor-only
diagnostic uses the same five-hour protocol on
typical failure case, with the final running-best conflict
values summarized by their range, mean, and median.

For the cross-framework transfer experiment, RLACO$_{\rm soft}$,
RLACO$_0$, RLACO$_{\rm config}$, and Pure ACO are evaluated at the
fixed target color budgets listed in Table~\ref{tab:rlaco-config}.
\texttt{flat1000\_76\_0} is excluded from the transfer experiment because its 1000-vertex size makes each ACO construction substantially more
expensive, resulting in too few completed ACO calls within a practical
wall-clock budget for a stable fixed-$k$ transfer comparison. Each instance uses an instance-specific wall-clock budget determined from preliminary RLACO runs to provide sufficient time for the running-best conflict value to stabilize. Stabilization is assessed by the relative improvement over the final 10% of the allocated budget, which is negligible across the tested instances.  The same instance-specific budget is then applied to
RLACO$_{\rm soft}$, RLACO$_0$, Pure ACO, and RLACO$_{\rm config}$ on that
instance. The main metric is the final running-best conflict
value at the fixed target $k$, reported as the observed best and
mean$\pm$sample standard deviation over ten runs. For the primary
RLACO$_{\rm soft}$--RLACO$_0$ comparison, two-sided Welch
two-sample $t$-tests are applied to the ten final conflict values, with
Holm correction across the ten instances. The phase-contribution
analysis is computed from the same RLACO$_{\rm soft}$ runs using
Eq.~\eqref{eq:phase-share}.

Finally, the RLACO$\rightarrow$Tabucol experiment uses the same
instance-specific RLACO preparation budget as the corresponding
transfer experiment, followed by the same five-hour downstream
Tabucol run. It also uses the same evaluation  as the
source-framework experiment: $k_{\mathrm{best}}$ is the primary metric,
with success count and best/mean successful-run time used as secondary
measures at comparable color budgets. For
RLACO$\rightarrow$Tabucol, the reported solving time starts when the
prepared coloring is handed to Tabucol; the preceding RLACO
preparation time is not included in this downstream timing measure because the timing measure is intended to evaluate
downstream Tabucol refinement from the RLACO-produced initialization.
